\documentclass[11pt]{article}
\usepackage{fontspec}
\usepackage{amsmath,amssymb}
\setmainfont{DejaVu Serif}
\usepackage[margin=1in]{geometry}
\title{TSR iteration 1: literature, an ordered clause-query DAG,\\ and a continuation-incompatibility cut}
\author{Grok, identity j15marti\\ repository \texttt{j15marti/pnp}, branch \texttt{main}}
\date{3 October 2026 (PT)}
\begin{document}
\maketitle
% GROK TSR ITER1 :: 2026-10-03_0950_GROK_TSR_ITER1

P versus NP remains open.

This note is iteration~1 of a time--space--representation pass.
It does not edit root \texttt{main.tex}.
It does not add a workflow.
The load-bearing compiler arrow below stays \textbf{OPEN}.
A finite \(W_\star(T)\) experiment is \textbf{COMPUTATIONAL} and does not close P versus NP.

The candidate architecture under study, as proposed on 2026-10-03, is not ``minimize DAG size.''
The resources are dynamic time \(T\) (longest source-to-answer path, or charged operations), mutable working memory \(S\), and stored representation or control size \(R\).
Continuation nonmergeability: histories \(H_1,H_2\) are incompatible when some common legal continuation \(Z\) gives different values of the decision \(F_{n,m}\).
\(K\) pairwise incompatible histories at one cut force at least \(K\) distinct states there, hence about \(\log_2 K\) bits besides the program counter.
Forgetting can be repaired by recomputation; pebbling is the proposed price of that repair, not a theorem of this note.
The open arrow is that an arbitrary SAT decider must compile into that charged model.
Status of the architecture: \textbf{CANDIDATE}.
Status of the compiler: \textbf{OPEN}.

\section{Part A. Literature actually opened}

Every bibliographic line below was read in this pass by WebSearch followed by WebFetch, or by opening the PDF text those tools returned.
Nothing in this section was copied from repository commentary.
Each item says what the opened text states, then what it does not say about \(D(n,m)\) or about P versus NP.

\subsection*{Beame, Jayram, and Saks}
Opened: ECCC TR98-053, authors Paul Beame, Michael Saks, and Jayram S.\ Thathachar, publication stamp 31 August 1998, including the abstract and the note that a corrected version was accepted on 17 September 1998.
Also opened: Paul Beame's time-space project page, which lists ``Paul Beame, T.\ S.\ Jayram, and Michael Saks, Time-space tradeoffs for branching programs, \emph{Journal of Computer and System Sciences} 63(4):542--572, December 2001,'' as a special issue of selected papers from FOCS 1998, and the FOCS 1998 item with the same three authors, pages 254--263.
The ECCC abstract states that there is a Boolean function \(f:\{0,1\}^n\to\{0,1\}\) that requires exponential size on every branching program of length \(cn\) for some constant \(c>1\), and that the paper separates syntactic and semantic read-\(k\) branching programs for \(k>1\) and gives an analogous length-versus-size tradeoff for \(R\)-way branching programs.
The journal article body was not opened, so this note does not import a theorem statement from JCSS.
The opened abstract does not mention \(D(n,m)\), ordered clause-query DAGs, or a proof that P differs from NP.
It is a \textbf{KNOWN THEOREM} about branching-program length and size for a constructed Boolean function, not a SAT time-space theorem.

\subsection*{Wegener}
Opened: Ingo Wegener, \emph{The Complexity of Boolean Functions}, Wiley--Teubner, copyright page 1987, ISBN 0\,471\,91555\,6, the LaTeX facsimile whose PDF metadata says the numbered items sit on the original pages.
Book page 418, immediately before \S14.2, states that the parity function has width-2 branching programs of depth \(n\), by testing \(x_{i+1}\) on level \(i\) and keeping the running parity in two nodes.
Section 14.2 defines \(f\) to be elusive when its branching-program depth equals \(n\), and Theorem 2.1 states \(\mathrm{BPD}(f)\ge l_{\max}(f)\ge c(f)\) for every \(f\in B_n\).
The proof shows that a depth-optimal path shorter than \(l_{\max}(f)\) would leave \(f\) constant on a large subcube.
That is a \textbf{KNOWN THEOREM} about branching-program depth.
It does not define \(D(n,m)\), does not price mutable memory against recomputation for SAT, and does not separate P from NP.

\subsection*{Fortnow}
Opened: Lance Fortnow, ``Time-Space Tradeoffs for Satisfiability,'' the manuscript at \texttt{lance.fortnow.com/papers/files/npvsnl.pdf}.
The manuscript does not print a journal year; a year is not assigned here.
Theorem 3.1 states that for every \(\varepsilon>0\),
\[
\mathrm{SAT}\notin\mathrm{NTISP}[n^{1+o(1)},n^{1-\varepsilon}].
\]
Corollary 3.3 states the same with \(\mathrm{DTISP}\) in place of \(\mathrm{NTISP}\).
The abstract adds that SAT cannot be solved deterministically in quasilinear time and \(\sqrt{n}\) space on the machines considered.
The paper says these bounds are likely far from the exponential time one expects, and it discusses an approach to NL versus NP together with relativized limits of that approach.
This is a \textbf{KNOWN THEOREM} for nondeterministic and deterministic time-space classes.
It does not mention \(D(n,m)\) and does not put SAT outside P.

\subsection*{Fortnow, Lipton, van Melkebeek, and Viglas}
Opened: Dieter van Melkebeek's page \texttt{pages.cs.wisc.edu/\textasciitilde dieter/Research/sat-ts.html}.
It names L.~Fortnow, R.~Lipton, D.~van Melkebeek, and A.~Viglas, \emph{Journal of the ACM} 52:835--865, 2005, and quotes the abstract.
The abstract states that for every constant \(c\) less than the golden ratio there is a positive constant \(d\) such that no deterministic random-access Turing machine solves satisfiability in time \(n^c\) and space \(n^d\), with \(d\) approaching 1 as \(c\) does, and that the same holds for conondeterministic machines for every constant \(c<\sqrt{2}\).
The proofs are by indirect diagonalization.
The abstract does not mention \(D(n,m)\).
The exponents are below 2, so the theorem does not separate P from NP.
Status: \textbf{KNOWN THEOREM}.

\subsection*{Williams and the exponent \(2\cos(\pi/7)\)}
Opened: Ryan Williams, ``Alternation-Trading Proofs, Linear Programming, and Lower Bounds (Extended Abstract),'' arXiv:1001.0746, as the ar5iv HTML.
The page header shows proceedings pagination 669--680 at Nancy.
The introduction states that the best known time lower bound for SAT on \(n^{o(1)}\)-space algorithms is \(n^{2\cos(\pi/7)-o(1)}\ge n^{1.801}\), obtained by a resource-trading proof cited there as Wil08, and that pushing the exponent to every constant would separate logarithmic space from NP.
Theorem 3.9 of the opened text states that if \(t\) and \(s\) are polynomially bounded, then any algorithm solving SAT in time \(t\) and space \(s\) requires \(t\cdot s=\Omega(n^{2\cos(\pi/7)-\varepsilon})\) for every \(\varepsilon>0\).
Theorem 3.8 states that there is no alternation-trading proof of \(\mathrm{NTIME}[n]\nsubseteq\mathrm{DTS}[n^2]\).
The numerical comparison \(1.801\) is the paper's, not a value computed in this run.
The paper does not mention \(D(n,m)\).
It does not prove a quadratic time lower bound, and it does not separate P from NP.
Status: \textbf{KNOWN THEOREM} inside the alternation-trading framework, with an explicit limitation of that framework.

\subsection*{Nordstr\"om}
Opened: Jakob Nordstr\"om, ``Pebble Games, Proof Complexity, and Time-Space Trade-offs,'' \emph{Logical Methods in Computer Science}, Volume 9, Issue 3, 13 September 2013, \texttt{doi:10.2168/LMCS-9(3:15)2013}, and the arXiv HTML of 1307.3913.
The abstract states that this is a survey of proof-complexity research that uses pebbling, focused on proof-space lower bounds and tradeoffs between proof size and proof space.
The opened body explains pebbling contradictions as CNF encodings of pebble games and surveys length-space tradeoffs for resolution and \(k\)-DNF resolution.
That is \textbf{KNOWN THEOREM} material about proof systems, and a structural analogy for ``recomputation has a pebbling price.''
It is not a theorem that an arbitrary SAT algorithm compiles into a charged multisign DAG, and it does not define or bound \(D(n,m)\).
It does not decide P versus NP.
Transfer class for this project: analogy requiring a new bridge, not an exact object transfer.

\subsection*{Atserias and Dalmau}
Opened: ECCC TR02-035, Albert Atserias and V\'ictor Dalmau, ``A Combinatorial Characterization of Resolution Width,'' ECCC publication stamp 18 June 2002, and the PDF served from \texttt{cs.upc.edu/\textasciitilde atserias/papers/width/ccc03.pdf}.
Both abstracts state that resolution width is characterized by the existential pebble game, and that the minimum space of refuting a 3-CNF is always at least the minimum width of refuting it, minus 3.
That is a \textbf{KNOWN THEOREM} relating resolution clause space and resolution width.
It does not lower-bound the state of an arbitrary SAT decider, does not mention \(D(n,m)\), and does not separate P from NP.
Same transfer class as the Nordstr\"om survey: proof-complexity analogy, missing the compiler theorem.

\subsection*{Dinesh, Koroth, and Sarma}
Opened: Krishnamoorthy Dinesh, Sajin Koroth, and Jayalal Sarma, ``Characterization and Lower Bounds for Branching Program Size using Projective Dimension,'' arXiv:1604.07200, the ar5iv HTML.
The abstract and Theorem 1.2 state that bitwise decomposable projective dimension characterizes deterministic branching-program size up to a polynomial: there is an absolute constant \(c>0\) such that
\[
\mathrm{bpdim}(f)/6\le\mathrm{bpsize}(f)\le(\mathrm{bpdim}(f))^c
\]
for every Boolean function \(f\), in the notation of that abstract.
Theorem 1.3 states that \(\mathrm{bitpdim}(\mathrm{SI}_d)\) is at least \(\Omega(d^3/\log d)\).
The introduction states that the best general explicit branching-program size lower bound remains Nechiporuk's \(\Omega(n^2/\log^2 n)\) bound, and that separating P from L would require exponential branching-program size for an explicit family.
The paper does not mention \(D(n,m)\) and does not separate P from NP, nor even P from L.
Status: \textbf{KNOWN THEOREM} about a linear-algebraic size measure.

\subsection*{What the literature licenses}
It licenses concrete time-space lower bounds for SAT on random-access machines inside indirect diagonalization, up through the alternation-trading exponent \(2\cos(\pi/7)\) and no farther by those proofs; branching-program length-versus-size and depth-versus-width facts; and resolution space-versus-width and pebbling translations inside proof complexity.
It does not license a Pareto theorem for \((T,S,R)\) on arbitrary SAT algorithms.
It does not license the compiler arrow.
That arrow remains \textbf{OPEN}.

\section{Part B. Build}

\subsection*{Corpus used for the definition}
Surveyed the tree of \texttt{j15marti/pnp} at commit \texttt{b75a119e11f4470a847dd05288846c159d55d31e}.
Read \texttt{AGENTS.md} (status vocabulary only; root \texttt{main.tex} was not edited, by the instruction for this pass).
Read \texttt{experiments/universal-dag/d5\_m8\_exact.cpp}, \texttt{research/explorations/random-challenge/2026-10-02\_1405\_LSPV5\_UNIVERSALDAG2.tex}, \texttt{results/raw/universal-dag/2026-10-02\_LSPV5\_UNIVERSALDAG2.txt}, \texttt{research/explorations/random-challenge/2026-10-02\_LSPV5\_FIXED\_NM\_DAG1.md}, and \texttt{results/raw/universal-dag/2026-10-02\_LSPV5\_UNIVERSAL3SAT1.txt}.
Did not open or write \texttt{research/agents/grok/2026-10-03\_0948\_GROK\_TSR1.tex}.

\subsection*{Definition that was coded}
Simple genuine 3-clauses: three distinct variables, one sign pattern in \(\{0,1\}^3\).
For \(n\) variables there are \(\binom{n}{3}\cdot 8\) clauses.
Each clause falsifies the assignments that match the pattern on its three variables and are free elsewhere, so it covers \(2^{n-3}\) assignments.
\(F_{n,m}=1\) when the selected set has size exactly \(m\) and the union of falsified assignments is the full cube; otherwise, on promise inputs, \(F_{n,m}=0\).
The input is the clause-presence vector in a fixed order, under the promise that exactly \(m\) bits are 1.
Dynamic state is \((t,k,M)\): clauses read, clauses selected, union of falsified assignments.
This is the state named in the UNIVERSALDAG2 log.

The reduction coded is the one in \texttt{d5\_m8\_exact.cpp}, including its prunes.
Terminals are \(0\) (promise-satisfiable), \(1\) (already unsatisfiable), and \(2\) (off-promise).
A state whose two children have the same identifier is not a decision node.
Two decision nodes at the same \(t\) with the same ordered pair of child identifiers are merged.
If the remaining clauses must all be taken to meet the weight promise, the C++ returns the determined terminal and does not allocate a node.
The sign-pattern round-robin order is: pattern \(0,\ldots,7\) outer, variable triples in lexicographic order inner.
The grouped order used as a check is: triples in lexicographic order, then the pattern.

Search limits: exact recursion with memoization from the root \((0,0,0)\).
No sampling, no order search, no symmetry quotient.
One Python~3 process on this computer.
The script that produced the first sweep is \texttt{/workspace/tsr\_iter1/ord\_dag.py}; the pairwise certificate was a second Python~3 process with the same recurrence.

\subsection*{What \(2107\) is, and what this run computed}
The repository log \texttt{2026-10-02\_LSPV5\_UNIVERSALDAG2.txt} states that a fixed lexicographic order has \(8\) nonterminal nodes for \(n=3,m=8\) and \(2107\) nonterminal nodes for \(n=4,m=8\), and it says \(2107\) is exact for that order only, not a minimum.
The log and the UNIVERSALDAG2 note both print \(2107\). No program source opened in this pass emits that count, and neither file records the lexicographic encoding.
This run did not reproduce \(2107\) under sign-tuple lexicographic order, mask lexicographic order, triple-then-pattern order, or sign round-robin, with or without the C++ tail collapse.
Therefore \(2107\) is not a number computed in this run, and it is not cited as one.
The \(n=3\) count \(8\) was reproduced only in the variant that does \emph{not} collapse a forced include-tail: that variant builds a single path of \(8\) decision nodes, width \(1\) at every level.
The C++ collapse, on the same \(n=3\) instance, returns the terminal \(1\) at the root and allocates \(0\) nodes, because the unique weight-\(8\) input is forced.
Both behaviors were computed.
They are different accountings of forced queries, not a disagreement about the one unsatisfiable formula.

Under the C++ accounting, which is the source semantics, this run computed the following for \(n=4\), \(m=8\).
\[
\binom{32}{8}=10518300,\qquad \text{UNSAT count}=272.
\]
Both match the repository logs, and both were recomputed here, not copied into the arithmetic.
Sign-pattern round-robin:
\[
\text{nonterminal nodes}=720,\quad
\text{memoized states}=1436,\quad
\text{longest decision path }T=31,
\]
\[
\text{level-}16\text{ width}=66
\quad(k\in\{2,3,4,5,6\}\text{ contributes }1+16+32+16+1).
\]
Triple-then-pattern order, same accounting: \(1890\) nonterminal nodes, peak width \(109\) at level \(14\).
The \(720\) and \(1890\) figures match the sign-round-robin and grouped columns of \texttt{FIXED\_NM\_DAG1.md} for \(m=8\).
That match is a re-verification.
The Gray column \(1094\) in that table was not reproduced; the binary-reflected pattern order tried here gave a different count, so \(1094\) is not reported as a result of this run.

The \(720\)-node program is an exact reduced ordered read-once DAG for one declared order (\textbf{PROVED/EXACT} for that order) and only an upper bound on the minimum over orders (\textbf{COMPUTATIONAL} as a minimization claim; the minimum remains \textbf{OPEN}).
It is not an FBDD minimum, not a circuit size, and not a lower bound on arbitrary algorithms.

\subsection*{Incompatibility lower bound, counted}
Fix the sign-round-robin program.
At cut \(t=16\) and selected-count \(k=4\), the reduction has \(32\) distinct continuation identifiers.
For every one of the \(\binom{32}{2}=496\) pairs, a search in the reduced DAG found a path on which the two identifiers reach opposite terminals in \(\{0,1\}\).
Paths that end in the off-promise terminal \(2\) were not accepted as witnesses.
Those \(32\) histories are pairwise incompatible in the sense of this note: some common continuation makes \(F_{4,8}\) differ.
Therefore any correct state at this fixed program counter and this fixed \(k\) needs at least \(32\) values.
The exhibited program has exactly \(32\) decision nodes in that cell, so the cell is tight.
\[
K=32,\qquad \log_2 K=5.
\]
The same-cut maximum of this same-\(k\) count, over all \((t,k)\), is this \(32\).
No larger same-\(k\) clique appeared in the memoized reachable states.

\section{Part C. Test}

Compare the counted lower bound with the exhibited program.

Inside the cell \((t,k)=(16,4)\), continuation incompatibility predicts \(32\) states and the program has \(32\) states.
On that slice the prediction meets the frontier.
Status of the comparison: \textbf{COMPUTATIONAL}.

It does not meet the full level width.
Level \(16\) has \(66\) decision nodes because five values of \(k\) are live.
Two histories with different \(k\) have no common legal continuation of weight exactly \(8\), since a shared bit-string preserves the difference in the number of clauses already chosen.
The incompatibility predicate therefore does not force those \(k\)-buckets apart.
The extra width is the selected-count, which is mutable memory, not a multisign separation.
So multisign incompatibility predicted the frontier of one fixed-\(k\) slice and did not predict the level width \(66\).
The \(n=3\) chain, width \(1\) and \(K=1\), is too small to tell: the predicate only recovers the trivial bound of one state.

Read as a Pareto point of the exhibited program, not as an optimum:
\[
T=31
\]
(longest source-to-answer chain of decision nodes),
\[
S\ge 5
\]
bits are forced inside the tight cell, while distinguishing all \(66\) nodes at level \(16\) takes \(\log_2 66\) bits of node identity plus the program counter (this run's \texttt{math.log2(66)} printed \(6.044394119358453\)),
\[
R=720
\]
stored decision nodes in this order, each with two children.
No claim is made that another order, an adaptive query policy, or a circuit cannot improve \(R\) or \(T\).
Minimizing \(R\) alone was not the objective.

The general \((T,S,R)\) tradeoff for arbitrary SAT algorithms is \textbf{OPEN}.
The compiler from an arbitrary algorithm into the charged model is \textbf{OPEN}.
One open load-bearing arrow means P versus NP is not proved.

\section{Status ledger}
\begin{itemize}
\item P versus NP: \textbf{OPEN}.
\item Beame--Jayram--Saks length-versus-size tradeoff, as in the opened ECCC abstract; Wegener Theorem 2.1; Fortnow's \(\mathrm{NTISP}/\mathrm{DTISP}\) bounds in the opened manuscript; the 2005 JACM golden-ratio and \(\sqrt{2}\) bounds; Williams's alternation-trading product bound and the stated \(2\cos(\pi/7)\) time exponent; Nordstr\"om's survey scope; Atserias--Dalmau space at least width minus 3; Dinesh--Koroth--Sarma bitwise projective dimension versus branching-program size: \textbf{KNOWN THEOREM}.
\item Sign-round-robin reduced DAG on \(D(4,8)\)'s clause-presence instance, C++ semantics: \(720\) nodes, \(T=31\), level-\(16\) width \(66\), UNSAT count \(272\), \(\binom{32}{8}=10518300\): \textbf{PROVED/EXACT} for that order and that accounting.
\item Same-cell incompatibility certificate \(K=32\) at \((t,k)=(16,4)\), all \(496\) pairs separated, matching the program on that cell: \textbf{COMPUTATIONAL} (finite, machine-checked).
\item \(D_{\mathrm{ord}}(4,8)=2107\): stated in a repository log for an underspecified lexicographic encoding; \emph{not reproduced this pass}; not used as a computed result.
\item Minimum over orders, adaptive DAGs, and circuits: \textbf{OPEN}.
\item Candidate \((T,S,R)\) architecture and pebbling-as-price: \textbf{CANDIDATE}.
\item Compiler from an arbitrary SAT decider: \textbf{OPEN}.
\item Nothing in this pass was \textbf{FALSIFIED/RETIRED}. The \(n=3\) collapse to \(0\) nodes shows only that charging forced queries changes \(R\); it does not retire either accounting.
\end{itemize}

\section{Next iteration}
The smallest honest step is not a larger \(m\) and not another order sweep.
It is to separate three storages that this run still folds together: the program counter \(t\), the selected-count \(k\), and the continuation identifier of \(M\).
On this same \(n=4,m=8\) sign-round-robin instance, charge a bit of \(k\) only when two different counts are simultaneously live, and charge a bit of the mask identifier only inside one count, using the \(K=32\) certificate already checked.
Then exhibit one recomputation step that drops \(k\) or \(M\) and pays for it by extra length, and record the resulting \((T,S,R)\) point as \textbf{COMPUTATIONAL}.
Do not treat that point as a lower bound on arbitrary algorithms until the compiler arrow is a theorem.

% END_OF_VERIFIED_REPORT::GROK::2026-10-03_0950_GROK_TSR_ITER1
\end{document}
